\documentclass[11pt]{article}

\usepackage{amsmath,amssymb}
\usepackage{xurl}
\usepackage[colorlinks=true,linkcolor=blue,citecolor=blue,urlcolor=blue]{hyperref}

\input{command}

\title{Wolf Corollary 6.7: Source-General Resolution Record}
\author{}
\date{2026-08-24}

\begin{document}
\maketitle
\gapnote{scope-restriction}{resolved}

\section{The assertion}

Corollary 6.7 of \emph{Quantum Channels \& Operations} (Wolf 2012) states: for a
trace-preserving, positive, linear map $T$ on $M_d(\mathbb{C})$ whose adjoint satisfies
the Schwarz inequality, and for any density matrix $\rho$ which is a \emph{maximum-rank}
fixed point of $T$, the set
\[
    \rho^{-1/2}\,\mathcal{F}_T\,\rho^{-1/2},
    \qquad \mathcal{F}_T = \{X \mid T(X) = X\},
\]
is a $*$-algebra, with the inverse square root taken on the support of $\rho$.

The corollary rests on the maximal-support property of fixed points stated earlier in
Section~6.4 of the reference: every fixed point of $T$ has support and range within the
support of the maximum-rank fixed point, so conjugation by $\rho^{-1/2}$ loses nothing.

\section{The formalization}

The source-general implementation is in
\leanid{QICLean/Channel/FixedPoint/AbstractWeightedFixedPoints.lean}.  It assumes exactly
that $T$ is positive and trace preserving and that its trace adjoint is Schwarz; it does
not assume a Kraus representation or complete positivity.

For an arbitrary positive semidefinite fixed point $\rho$ with support projection $Q$,
\leanid{IsPositiveMap.weightedCornerFixedPointsStarSubalgebra} proves that
\[
    \left\{Y\in QM_D(\mathbb C)Q\mid
      T(\sqrt\rho\,Y\sqrt\rho)=\sqrt\rho\,Y\sqrt\rho\right\}
\]
is a $*$-subalgebra of $QM_D(\mathbb C)Q$.  Its multiplication proof follows Wolf's
Theorem~6.14 density blocks.  On the support sector, fixed blocks are
$\sigma_k\otimes X_k$ and $\rho$ has blocks $\sigma_k\otimes R_k$ with both density
factors positive definite.  The weighted product remains fixed because
\[
  (\sigma_k\otimes X_k)(\sigma_k^{-1}\otimes R_k^{-1})
  (\sigma_k\otimes Z_k)
  =\sigma_k\otimes(X_kR_k^{-1}Z_k).
\]
This is the factor-swapped form of Wolf's
$\mathcal M_{d_k}\otimes\rho_k$ notation, as already recorded for the formalized
Theorem~6.14.

The every-maximum-rank step is source-general as well.
\leanid{IsPositiveMap.maximalSupport\_of\_maximalRank} in
\leanid{QICLean/Channel/FixedPoint/AbstractMaximalRank.lean} takes an arbitrary positive
semidefinite fixed point whose rank bounds the rank of every stationary density and
proves $QXQ=X$ for every fixed point $X$.  It compares that arbitrary choice with
$T_\infty(\mathbf 1)$, rather than replacing Wolf's quantifier by the canonical witness.

Finally,
\leanid{IsPositiveMap.mem\_weightedCornerFixedPointsStarSubalgebra\_iff\_inverseSandwich}
identifies the formal carrier in both directions with
$\rho^{-1/2}\mathcal F_T\rho^{-1/2}$.  The inverse square root is the functional-calculus
inverse on $\operatorname{supp}(\rho)$ and zero on its kernel.  The support identities
$\sqrt\rho\,\rho^{-1/2}=Q=\rho^{-1/2}\sqrt\rho$ give the equality without an ambient
invertibility assumption.  The generic construction also includes zero support as the
zero corner algebra; the source theorem's density hypothesis, of course, has trace one.

\section{Verdict}

The former Kraus-only scope restriction is eliminated.
\leanid{IsPositiveMap.wolfCorollary67} states the printed corollary for a positive
trace-preserving linear map with Schwarz trace adjoint and for every maximum-rank
stationary density matrix, and it returns a star-algebra whose carrier is exactly Wolf's
displayed inverse-square-root set.  The older declarations in the
\texttt{Kraus} namespace remain compatibility specializations; the maximum-rank and
surjectivity declarations delegate to the source-general proofs rather than furnishing
a parallel route.

\end{document}
